\documentclass[openany]{book}

\usepackage[margin=1in]{geometry}
\usepackage{amsmath,amsfonts,amsthm, amssymb}
\usepackage{yhmath}
\usepackage{mathrsfs}
\usepackage{mathtools}
\usepackage{xcolor}
\usepackage{graphicx}
\usepackage{comment}
\usepackage{tikz-cd}
\usepackage{quiver}
\renewcommand{\familydefault}{ppl}
\newcommand{\tr}{\text{tr}}
\newcommand{\R}{\mathbb{R}}
\newcommand{\E}{\mathbb{E}}
\newcommand{\Z}{\mathbb{Z}}
\newcommand{\C}{\mathbb{C}}
\newcommand{\F}{\mathbb{F}}
\newcommand{\la}{\langle}
\newcommand{\ra}{\rangle}
\newcommand{\colim}{\text{colim}}
\DeclareMathOperator{\im}{im}
\let\oldemptyset\emptyset
\let\emptyset\varnothing
\newcommand{\tor}{\text{Tor}}
\newcommand{\id}{\text{id}}
\newcommand{\ext}{\text{Ext}}
\newcommand{\ptop}{\text{PTop}}
\newcommand{\pt}{\text{pt}}
\newcommand{\ach}{\text{Ach}}
\newcommand{\Q}{\mathbb{Q}}
\newcommand{\gal}{\text{Gal}}
\newcommand{\diverg}{\operatorname{div}}
\newcommand{\curl}{\operatorname{curl}}

\usepackage{thmtools,thm-restate}

% Fixing mdframed skip below
% See https://tex.stackexchange.com/a/292090/143086
\usepackage[framemethod=TikZ]{mdframed}
\usepackage{xpatch}
\makeatletter
\xpatchcmd{\endmdframed}
	{\aftergroup\endmdf@trivlist\color@endgroup}
	{\endmdf@trivlist\color@endgroup\@doendpe}
	{}{}
\makeatother

\definecolor{huilightpink}{HTML}{fcf2f9}
\definecolor{huidarkpink}{HTML}{ed34b3}
\declaretheoremstyle[
	mdframed={
		backgroundcolor=huilightpink,
		linecolor=huidarkpink,
		rightline=false,
		topline=false,
		bottomline=false,
		linewidth=2pt,
		innertopmargin=5pt,
		innerbottommargin=8pt,
		innerleftmargin=8pt,
		leftmargin=-2pt,
		skipbelow=2pt,
		nobreak
	},
	headfont=\normalfont\bfseries\color{huidarkpink}
]{huipinkbox}
\declaretheorem[style=huipinkbox,name=Theorem,within=chapter]{thm}
\declaretheorem[style=huipinkbox,name=Theorem,sibling=thm]{theorem}





\definecolor{huilightyellow}{HTML}{fff5d6}
\definecolor{huidarkyellow}{HTML}{fcad03}
\declaretheoremstyle[
	mdframed={
		backgroundcolor=huilightyellow,
		linecolor=huidarkyellow,
		rightline=false,
		topline=false,
		bottomline=false,
		linewidth=2pt,
		innertopmargin=5pt,
		innerbottommargin=8pt,
		innerleftmargin=8pt,
		leftmargin=-2pt,
		skipbelow=2pt,
		nobreak
	},
	headfont=\normalfont\bfseries\color{huidarkyellow}
]{huiyellowbox}
\declaretheorem[style=huiyellowbox,name=Proposition,within=chapter]{prop}

\definecolor{huilightpurple}{HTML}{faf2ff}
\definecolor{huidarkpurple}{HTML}{912ed9}
\declaretheoremstyle[
	mdframed={
		backgroundcolor=huilightpurple,
		linecolor=huidarkpurple,
		rightline=false,
		topline=false,
		bottomline=false,
		linewidth=2pt,
		innertopmargin=5pt,
		innerbottommargin=8pt,
		innerleftmargin=8pt,
		leftmargin=-2pt,
		skipbelow=2pt,
		nobreak
	},
	headfont=\normalfont\bfseries\color{huidarkpurple}
]{huipurplebox}
\declaretheorem[style=huipurplebox,name=Lemma,within=chapter]{lem}


\definecolor{huilightpurple}{HTML}{faf2ff}
\definecolor{huidarkpurple}{HTML}{912ed9}
\declaretheoremstyle[
	mdframed={
		backgroundcolor=huilightpurple,
		linecolor=huidarkpurple,
		rightline=false,
		topline=false,
		bottomline=false,
		linewidth=2pt,
		innertopmargin=5pt,
		innerbottommargin=8pt,
		innerleftmargin=8pt,
		leftmargin=-2pt,
		skipbelow=2pt,
		nobreak
	},
	headfont=\normalfont\bfseries\color{huidarkpurple}
]{huipurplebox}
\declaretheorem[style=huipurplebox,name=Definition,within=chapter]{defn}

\definecolor{huilightblue}{HTML}{edf9ff}
\definecolor{huidarkblue}{HTML}{4b79db}
\declaretheoremstyle[
	mdframed={
		backgroundcolor=huilightblue,
		linecolor=huidarkblue,
		rightline=false,
		topline=false,
		bottomline=false,
		linewidth=2pt,
		innertopmargin=5pt,
		innerbottommargin=8pt,
		innerleftmargin=8pt,
		leftmargin=-2pt,
		skipbelow=2pt,
		nobreak
	},
	headfont=\normalfont\bfseries\color{huidarkblue}
]{huiblueblox}
\declaretheorem[style=huiblueblox,name=Example,within=chapter]{example}

\declaretheoremstyle[
	mdframed={
		backgroundcolor=huilightblue,
		linecolor=huidarkblue,
		rightline=false,
		topline=false,
		bottomline=false,
		linewidth=2pt,
		innertopmargin=5pt,
		innerbottommargin=8pt,
		innerleftmargin=8pt,
		leftmargin=-2pt,
		skipbelow=2pt,
		nobreak
	},
	headfont=\normalfont\bfseries\color{huidarkblue}
]{huiblueblox}
\declaretheorem[style=huiblueblox,name=Problem,within=chapter]{prob}

\newcommand{\nirwarnsymbol}{%
	\begin{tikzpicture}[baseline=(x.base)]
		\draw[rounded corners=.01em] (-.05em,-1.07em)rectangle(.05em,.78em);
		\draw[fill=white,rounded corners=1.3] (0,.75em)--(.75em,0)--(0,-.75em)--(-.75em,0)--cycle;
		\draw[line width=0.2mm, line cap=round](-.4em,-1.07em)--(.4em,-1.07em);
		\node(x) at (0,0em) {};
		% Thank you https://tex.stackexchange.com/a/262510
		\draw[
			line cap=but,
			line join=round,
			x=.5em,
			line width=0.5mm,
			y=1*(height("Z")-\pgflinewidth)*(1-sin(10)),
			rotate=-10,
			rounded corners=1.5pt,
		](-0.57, 0.57) -- (0.57, 0.57) -- (-0.57, -0.57) -- (0.57, -0.57);
	\end{tikzpicture}%
}

%%%%%%%%%%%%%%%%%%%%%%%%%%%%%%%%%%%%%%%%%%%% MARGINS
\usepackage{marginnote}
% Thank you https://tex.stackexchange.com/a/472882
% Makes marginnotes always appear on the left, apparently
%
\makeatletter
\long\def\@mn@@@marginnote[#1]#2[#3]{%
	\begingroup
		\ifmmode\mn@strut\let\@tempa\mn@vadjust\else
			\if@inlabel\leavevmode\fi
			\ifhmode\mn@strut\let\@tempa\mn@vadjust\else\let\@tempa\mn@vlap\fi
		\fi
		\@tempa{%
			\vbox to\z@{%
				\vss
				\@mn@margintest
				\if@reversemargin\if@tempswa
						\@tempswafalse
					\else
						\@tempswatrue
				\fi\fi

					\llap{%
						\vbox to\z@{\kern\marginnotevadjust\kern #3
							\vbox to\z@{%
								\hsize\marginparwidth
								\linewidth\hsize
								\kern-\parskip
								%\mn@parboxrestore
								\marginfont\raggedleftmarginnote\strut\hspace{\z@}%
								\ignorespaces#1\endgraf
								\vss
							}%
							\vss
						}%
						\if@mn@verbose
							\PackageInfo{marginnote}{xpos seems to be \@mn@currxpos}%
						\fi
						\begingroup
							\ifx\@mn@currxpos\relax\else\ifx\@mn@currpos\@empty\else
									\kern\@mn@currxpos
							\fi\fi
							\ifx\@mn@currpage\relax
								\let\@mn@currpage\@ne
							\fi
							\if@twoside\ifodd\@mn@currpage\relax
									\kern-\oddsidemargin
								\else
									\kern-\evensidemargin
								\fi
							\else
								\kern-\oddsidemargin
							\fi
							\kern-1in
						\endgroup
						\kern\marginparsep
					}%
			}%
		}%
	\endgroup
}
\makeatother
%
% Mostly for todonotes
\renewcommand{\marginpar}{\marginnote}
%%%%%%%%%%%%%%%%%%%%%%%%%%%%%%%%%%%%%%%%%%%% /MARGINS

\definecolor{nirlightred}{RGB}{250, 220, 220}
\definecolor{nirdarkred}{HTML}{f40000}
\declaretheoremstyle[
	mdframed={
		backgroundcolor=nirlightred,
		linecolor=nirdarkred,
		rightline=false,
		topline=false,
		bottomline=false,
		linewidth=2pt,
		innertopmargin=5pt,
		innerbottommargin=8pt,
		innerleftmargin=8pt,
		leftmargin=-2pt,
		skipbelow=2pt,
		nobreak
	},
	headfont=\normalfont\bfseries\color{nirdarkred}
]{nirredbox}

% \makeatletter
% \declaretheorem[
% 	style=nirredbox,
% 	name=Warning,
% 	sibling=thm,
% 	% without \leavevmode, the first item in a list gets misformatted
% 	postheadhook={\leavevmode\marginnote{\nirwarnsymbol}[-3pt]%
% 	\ifthmt@thisistheone% restatable makes alignment weird
% 		\hspace{-2.2pt}%
% 	\fi}
% ]{warn}
% \makeatother

\newcommand{\nirideasymbol}{%
	\begin{tikzpicture}[baseline=(x.base)]
		\draw[rounded corners=.01em] (-.05em,-1.07em)rectangle(.05em,.78em);
		\draw[fill=white,rounded corners=1.3] (0,.75em)--(.75em,0)--(0,-.75em)--(-.75em,0)--cycle;
		\draw[line width=0.2mm, line cap=round](-.4em,-1.07em)--(.4em,-1.07em);
		\node(x) at (0,0em) {};
		\node at (0,0em) {{\textbf{!}}};
	\end{tikzpicture}%
}
\renewcommand{\nirwarnsymbol}{%
	\begin{tikzpicture}[baseline=(x.base)]
		\draw[rounded corners=.01em] (-.05em,-1.07em)rectangle(.05em,.78em);
		\draw[fill=white,rounded corners=1.3] (0,.75em)--(.75em,0)--(0,-.75em)--(-.75em,0)--cycle;
		\draw[line width=0.2mm, line cap=round](-.4em,-1.07em)--(.4em,-1.07em);
		\node(x) at (0,0em) {};
		% Thank you https://tex.stackexchange.com/a/262510
		\draw[
			line cap=but,
			line join=round,
			x=.5em,
			line width=0.5mm,
			y=1*(height("Z")-\pgflinewidth)*(1-sin(10)),
			rotate=-10,
			rounded corners=1.5pt,
		](-0.57, 0.57) -- (0.57, 0.57) -- (-0.57, -0.57) -- (0.57, -0.57);
	\end{tikzpicture}%
}
\makeatletter
\declaretheorem[
	style=nirredbox,
	name=Idea,
	sibling=thm,
	% without \leavevmode, the first item in a list gets misformatted
	postheadhook={\leavevmode\marginnote{\nirideasymbol}[-3pt]%
	\ifthmt@thisistheone% restatable makes alignment weird
		\hspace{-2.2pt}%
	\fi}
]{idea}

\declaretheorem[
	style=nirredbox,
	name=Warning,
	sibling=thm,
	% without \leavevmode, the first item in a list gets misformatted
	postheadhook={\leavevmode\marginnote{\nirwarnsymbol}[-3pt]%
	\ifthmt@thisistheone% restatable makes alignment weird
		\hspace{-2.2pt}%
	\fi}
]{warn}
\makeatother

\title{Calc III Sections
\\ 
\vspace{0.4cm}
\large Fall 2025}






\date{\today}
\author{Hui Sun}


\begin{document}

\maketitle

% \tableofcontents
\newpage


\begin{center}
    \Large Calc III-Week 11 (11/3-7)
\end{center}

\renewcommand\thesection{\arabic{section}}

\noindent
Topic: theorem of change of variables.




\begin{defn}[injective, surjective]
    Let $T:\R^2\to\R^2$ be a map, then $T$ is \textbf{injective or one-to-one} if for $x,y\in\R^2$
    \begin{equation*}
        Tx=Ty
    \end{equation*}
    then
    \begin{equation*}
        x=y.
    \end{equation*}
    We say $T$ is \textbf{surjective or onto} if for all $y\in\R^2$, there exists $x\in\R^2$ such that 
    \begin{equation*}
        Tx=y
    \end{equation*}
    If $T$ is both injective and bijective, then we say $T$ is \textbf{bijective}.
\end{defn}

\begin{defn}[linear map]
    Suppose $T:\R^2\to\R^2$, then $T$ is linear if and only if for all $x,y\in\R^2, \lambda\in\R$.
    \begin{equation*}
        \begin{cases}
            T(x+y)=Tx+Ty\\
            T(\lambda x)=\lambda Tx
        \end{cases}
    \end{equation*}
\end{defn}


\begin{defn}[Jacobian Determinant]
    Let $T:D^*\subset\R^2\to\R^2$ be of $C^1$ defined by 
    \begin{equation*}
        T:\begin{pmatrix}
            u\\
            v
        \end{pmatrix}\mapsto \begin{pmatrix}
            x(u,v)\\
            y(u,v)
        \end{pmatrix}
    \end{equation*}
    The \textbf{Jacobian determinant} of $T$, denoted as $\frac{\partial(x,y)}{\partial(u,v)}$ is the determinant of the matrix $DT(u,v)$:
    \begin{equation*}
        \frac{\partial(x,y)}{\partial(u,v)}=\det\begin{bmatrix}
            \frac{\partial x}{\partial u}& \frac{\partial x}{\partial v}\\
            \frac{\partial y}{\partial u}&\frac{\partial y}{\partial v}
        \end{bmatrix}
    \end{equation*}
\end{defn}


\begin{example}
    The polar coordinate transformation is given by:
    \begin{equation*}
        \begin{pmatrix}
            x\\
            y
        \end{pmatrix}\mapsto \begin{pmatrix}
            r\cos\theta\\
            r\sin\theta
        \end{pmatrix}
    \end{equation*}
    The Jacobian matrix is given by 
    \begin{equation*}
        \det\begin{bmatrix}
            \cos\theta&-r\sin\theta\\
            \sin\theta&r\cos\theta
        \end{bmatrix}=r
    \end{equation*}
\end{example}



\begin{prop}
    Let $T:\R^2\to\R^2$ be a linear map, then there exists a $2\times 2$ matrix $A=\begin{pmatrix}
        a_{11}&a_{12}\\
        a_{21}&a_{22}
    \end{pmatrix}$ such that 
    \begin{equation*}
        Tx=Ax
    \end{equation*}
\end{prop}


\begin{prop}
    Let $T:\R^2\to\R^2$ be a linear map, then $T$ is injective if and only if it is also surjective if and only if $\det(A)\neq 0$, where $A$ is the matrix associated to $T$.
\end{prop}


\begin{thm}[change of variables formula]
    Let $D,D^*$ be elementary regions in $\R^2$, suppose $T:D^*\to D$ is bijective. Then for any integral function $f:D\to\R$, the \textbf{change of variable formula} states
    \begin{equation*}
        \iint_Df(x,y)dxdy=\iint_{D^*}f(x(u,v),y(u,v))\left|\det(J)\right|dudv
    \end{equation*}
    where 
    \begin{equation*}
        \det(J)=\left|\frac{\partial(x,y)}{\partial(u,v)}\right|
    \end{equation*}
    is the determinant of the Jacobian matrix.
\end{thm}

\begin{cor}
    An immediate corollary of the above theorem is the change of variables into polar cooridnates:
    \begin{equation*}
        \iint_Df(x,y)dxdy=\iint_{D^*}f(r\cos\theta,r\sin\theta)rdrd\theta
    \end{equation*}
\end{cor}

\begin{cor}
    Let $W,W^*$ be elementary regions in $\R^3$, and suppose $T:W^*\to W$ is bijective. Then the change of varialbes formula for triple integrals states:
    \begin{equation*}
        \iiint_Wf(x,y,z)dxdydz=\iiint_{W^*}f(x(u,v,w), y(u,v,w), z(u,v,w))\mid\det(J)\mid dudvdw
    \end{equation*}
    where 
    \begin{equation*}
        \det(J)=\det\begin{bmatrix}
            \frac{\partial x}{\partial u}&\frac{\partial x}{\partial v}&\frac{\partial x}{\partial w}\\
            \frac{\partial y}{\partial u}&\frac{\partial y}{\partial v}&\frac{\partial y}{\partial w}\\
            \frac{\partial z}{\partial u}&\frac{\partial z}{\partial v}&\frac{\partial z}{\partial w}
        \end{bmatrix}
    \end{equation*}
\end{cor}


\begin{prob}
    Compute the following integral:
    \begin{equation*}
        \iint_D(x^2+y^2)^{1/2}dxdy
    \end{equation*}
    where $D$ is the disk $D=\{(x,y):x^2+y^2\leq 9\}$.
\end{prob}
\begin{proof}
    Writing this in polar coordinates and using the change of variables formula, we have 
    \begin{align*}
        \iint_D(x^2+y^2)^{1/2}dxdy&=\int_0^{2\pi}\int_{0}^3 (r^2)^{1/2}rdrd\theta\\
        &=\int_0^{2\pi}\int_0^3r^2drd\theta\\
        &=18\pi
    \end{align*}
\end{proof}

\begin{prob}[Marsden-Tromba, 6.2, Exercise 19]
    Compute $\iint_R(x+y)^2e^{x-y}dxdy$, where $R$ is the region bounded by $x+y=1, x+y=4, x-y=-1, x-y=1$.
\end{prob}
% \begin{proof}
%     Writing $u=x+y, v=x-y$, by the change of variables formula, we have 
%     \begin{align*}
%         \iint_R(x+y)^2e^{x-y}dxdy&=\int_{-1}^1\int_1^4u^2e^{v}|\det(J)|dudv\\
%         &=\frac{1}{2}\left(\int_1^4u^2du\right)\left(\int_{-1}^1e^vdv\right)\\
%         &=\frac{21}{2}\left(e-\frac{1}{e}\right)
%     \end{align*}
% \end{proof}




\begin{prob}
    Compute $\iint_RxydA$ where $R$ is the region bounded by $xy=1, xy=3, y=2, y=6$, using the transformation $x=\frac{v}{6u}, y=2u$.
\end{prob}
% \begin{proof}
%     \begin{align*}
%         \iint_Rxy^2dA&=\int_1^3\int_3^9\frac{v}{6u}4u^2|\det(J)|dvdu\\
%         &=\int_1^3\int_3^9\frac{2v}{9}dvdu\\
%         &=\frac{2}{3}\int_3^9vdv\\
%         &=24
%     \end{align*}
%     where 
%     \begin{equation*}
%         \det(J)=\det\begin{bmatrix}
%             -\frac{v}{6u^2}&\frac{1}{6u}\\
%             2&0
%         \end{bmatrix}=-\frac{1}{3u}
%     \end{equation*}
% \end{proof}










% change of variables examples (of double and triple integrals)


\end{document}